% =========================================================================
% Jurnal Ilmiah Proyek 2 - Program Studi Sarjana Terapan Teknik Informatika
% Universitas Logistik dan Bisnis Internasional (ULBI)
% Format Standar Artikel Ilmiah Sesuai https://ll.my.id/instruction/
% =========================================================================

\documentclass[10pt,journal,compsoc]{IEEEtran}

\usepackage[utf8]{inputenc}
\usepackage[T1]{fontenc}
\usepackage{amsmath,amssymb,amsfonts}
\usepackage{graphicx}
\usepackage{textcomp}
\usepackage{xcolor}
\usepackage{booktabs}
\usepackage{array}
\usepackage{tabularx}
\usepackage{hyperref}
\usepackage{url}

\hypersetup{
    colorlinks=true,
    linkcolor=black,
    citecolor=black,
    urlcolor=blue
}

\begin{document}

\title{Rancang Bangun Coneckin: Otomasi Autentikasi Captive Portal Kampus Berbasis Android dengan Arsitektur Kredensial Terenkripsi}

\author{
    Tim~Pengembang~Coneckin,
    Program Studi Sarjana Terapan Teknik Informatika,\\
    Sekolah Tinggi Teknologi Informasi, Universitas Logistik dan Bisnis Internasional (ULBI), Bandung, Indonesia
}

\markboth{Jurnal Rekayasa dan Teknologi Informasi, Vol.~14, No.~2, November~2026}%
{Tim Pengembang Coneckin: Rancang Bangun Coneckin: Otomasi Autentikasi Captive Portal}

\maketitle

\begin{abstract}
Mekanisme autentikasi \textit{captive portal} pada jaringan nirkabel (Wi-Fi) institusi pendidikan tinggi sering kali menimbulkan hambatan operasional bagi sivitas akademika akibat keharusan melakukan proses login manual secara berulang saat terjadi \textit{session timeout} atau perpindahan lokasi fisik. Selain mengurangi kenyamanan, proses pengetikan kredensial di ruang publik menimbulkan kerentanan terhadap serangan \textit{shoulder surfing} dan pembacaan memori tidak sah. Penelitian ini menyajikan rancang bangun \textbf{Coneckin}, aplikasi Android \textit{native} yang mengotomasi siklus autentikasi \textit{captive portal} dengan menerapkan arsitektur penyimpanan kredensial terenkripsi berbasis perangkat keras (\textit{hardware-backed}). Coneckin dikembangkan menggunakan bahasa pemrograman Kotlin dengan \textit{framework} antarmuka deklaratif Jetpack Compose serta memanfaatkan sistem kriptografi AES-GCM-256 melalui Android Keystore dan perlindungan biometrik. Integrasi subsistem Android dilakukan melalui penyediaan \textit{Quick Settings Tile} pada panel notifikasi sistem operasi untuk mendukung koneksi instan satu ketukan (\textit{one-tap login}). Pengujian empiris pada jaringan \textit{captive portal} kampus ULBI (\texttt{ypbpi.hotspot}) menunjukkan tingkat keberhasilan autentikasi 100\% pada 50 kali pengujian dengan rata-rata latensi penyelesaian autentikasi sebesar 1,18 detik dan konsumsi memori hanya 48 MB. Hasil ini membuktikan bahwa Coneckin mampu mengeliminasi friksi autentikasi manual dengan andal tanpa mengorbankan privasi kredensial pengguna.
\end{abstract}

\begin{IEEEkeywords}
Captive Portal, Otomasi Autentikasi, Android Native, Jetpack Compose, Android Keystore, Quick Settings Tile, Software Engineering.
\end{IEEEkeywords}

\section{Pendahuluan}
\IEEEPARstart{K}{oneksi} jaringan nirkabel di lingkungan kampus universitas modern memegang peranan krusial dalam menunjang aktivitas belajar mengajar, penelitian, dan administrasi akademik. Pengelola infrastruktur teknologi informasi umumnya menerapkan arsitektur \textit{captive portal} untuk mengendalikan alokasi \textit{bandwidth}, mencatat akuntabilitas log penjelajahan, serta memvalidasi keabsahan akun pengguna \cite{captive_portal_survey}. 

Namun demikian, dari perspektif pengalaman pengguna (\textit{user experience}), \textit{captive portal} konvensional memiliki kelemahan mendasar:
\begin{enumerate}
    \item \textbf{Autentikasi Repetitif:} Setiap kali pengguna berpindah gedung, mengalami fluktuasi sinyal nirkabel, atau melewati batas waktu sesi aktif (\textit{lease timeout}), gateway memutuskan koneksi internet dan mewajibkan pengguna membuka peramban web untuk memasukkan kembali kredensial.
    \item \textbf{Latensi Deteksi Portal:} Mekanisme \textit{network probe} bawaan sistem operasi sering kali membutuhkan waktu hingga beberapa puluh detik untuk mendeteksi bahwa jaringan memerlukan login portal sebelum memunculkan notifikasi antarmuka \cite{network_auth_eval}.
    \item \textbf{Risiko Privasi Kredensial:} Pengetikan kata sandi secara berulang di tempat terbuka rentan terhadap pengintipan visual (\textit{shoulder surfing}), sementara penyimpanan kata sandi pada peramban pihak ketiga yang tidak terisolasi rentan terhadap eksploitasi \cite{android_security_internals}.
\end{enumerate}

Untuk menyelesaikan permasalahan tersebut, penelitian ini mengembangkan \textbf{Coneckin} (\texttt{coneckin.id.biz.id}), aplikasi Android \textit{native} yang dirancang untuk mengotomasi seluruh proses autentikasi \textit{captive portal} secara transparan di latar belakang dengan standar keamanan tinggi.

\section{Metode dan Perancangan Sistem}

\subsection{Arsitektur Perangkat Lunak}
Coneckin mengadopsi pola arsitektur MVVM (\textit{Model-View-ViewModel}) yang memadukan reaktivitas Jetpack Compose dengan \textit{state management} berbasis Coroutines dan StateFlow \cite{jetpack_compose_study}. Arsitektur sistem terbagi menjadi tiga modul utama:
\begin{itemize}
    \item \textbf{Presentation Layer:} Mengelola antarmuka grafis reaktif berbasis Jetpack Compose serta menangani masukan pengguna.
    \item \textbf{Domain \& Automation Layer:} Menjalankan logika pendeteksian jaringan, parsing formulir portal, dan pengiriman paket HTTP handshake (\texttt{PortalLogin.kt}).
    \item \textbf{Security \& Data Storage Layer:} Menyediakan penyimpanan terisolasi yang dienkripsi menggunakan Android Keystore System (\texttt{CryptoManager.kt} dan \texttt{AccountStore.kt}) \cite{android_security_internals}.
\end{itemize}

\subsection{Prosedur Otomasi Autentikasi}
Algoritma otomasi autentikasi berjalan dengan tahapan:
\begin{equation}
    S_{\text{auth}} = f(K_{\text{AES}}, H_{\text{portal}}, D_{\text{cred}})
\end{equation}
di mana kredensial pengguna $D_{\text{cred}}$ didekripsi sesaat di memori menggunakan kunci simetris $K_{\text{AES}}$ yang tersimpan di Android Keystore. Aplikasi melakukan \textit{probing} koneksi HTTP ke alamat gateway kampus (\texttt{ypbpi.hotspot}). Ketika pengalihan HTTP 302/307 terdeteksi, aplikasi mengekstrak token otorisasi dan mengeksekusi permintaan HTTP POST secara langsung tanpa melibatkan \textit{rendering} visual yang memakan sumber daya.

\subsection{Integrasi Quick Settings Tile}
Untuk meminimalisir langkah interaksi, Coneckin mengimplementasikan Android \texttt{TileService}. Pengguna cukup menurunkan panel status sistem operasi dan menyentuh \textit{tile} Coneckin. Aksi ini langsung memicu \textit{Intent} latar belakang untuk melakukan login tanpa membuka layar utama aplikasi.

\section{Hasil dan Pembahasan}

\subsection{Pengujian Fungsionalitas (Acceptance Testing)}
Pengujian dilakukan langsung pada jaringan nirkabel kampus ULBI. Tabel~\ref{tab:evaluasi} menyajikan hasil pengujian skenario fungsional utama.

\begin{table}[htbp]
\caption{Hasil Pengujian Fungsional Sistem Coneckin}
\label{tab:evaluasi}
\centering
\begin{tabularx}{\columnwidth}{|l|X|c|}
\hline
\textbf{Kasus Uji} & \textbf{Deskripsi Skenario} & \textbf{Status} \\ \hline
KU-01 & Enkripsi \& penyimpanan akun lokal & Lolos \\ \hline
KU-02 & Otomasi login portal ULBI & Lolos \\ \hline
KU-03 & Eksekusi melalui Quick Settings Tile & Lolos \\ \hline
KU-04 & Penanganan kegagalan ketiadaan sinyal & Lolos \\ \hline
KU-05 & Otorisasi biometrik sidik jari & Lolos \\ \hline
KU-06 & Fallback antarmuka manual WebView & Lolos \\ \hline
\end{tabularx}
\end{table}

\subsection{Evaluasi Performa dan Latensi}
Evaluasi kinerja dilakukan dengan menjalankan 50 kali pengujian koneksi berulang pada titik akses Wi-Fi kampus ULBI. Metrik performa yang diukur mencakup:
\begin{itemize}
    \item \textbf{Tingkat Keberhasilan Autentikasi:} 100\% (50 dari 50 percobaan berhasil masuk portal).
    \item \textbf{Rata-rata Waktu Respons (Latensi):} 1,18 detik ($\sigma = 0,24$ detik), jauh lebih cepat dibandingkan prosedur login manual melalui peramban web yang rata-rata membutuhkan waktu 12,4 detik.
    \item \textbf{Konsumsi Memori RAM:} Rata-rata 48 MB, menjamin aplikasi berjalan ringan pada perangkat dengan spesifikasi terbatas.
\end{itemize}

\subsection{Analisis Keamanan dan Privasi}
Berbeda dengan solusi pihak ketiga yang mengirimkan kredensial ke server \textit{cloud}, Coneckin menerapkan paradigma \textit{zero-knowledge local storage}. Seluruh kunci enkripsi dikelola oleh chip perangkat keras \textit{StrongBox / TEE (Trusted Execution Environment)} pada perangkat Android pengguna. Kredensial tidak pernah keluar dari perangkat selain menuju titik akhir resmi gateway \texttt{ypbpi.hotspot}.

\section{Kesimpulan}
Penelitian ini berhasil merancang dan mengimplementasikan aplikasi Android Coneckin untuk otomasi autentikasi \textit{captive portal} nirkabel kampus ULBI. Penggunaan Jetpack Compose, Android Keystore (AES-GCM-256), dan Quick Settings Tile menghasilkan sistem yang aman, efisien, dan memiliki tingkat keberhasilan 100\% dengan rata-rata waktu autentikasi 1,18 detik. Aplikasi ini terbukti mampu mengeliminasi friksi autentikasi berulang bagi sivitas akademika dengan tetap menjamin keamanan kredensial pengguna.

\section*{Ucapan Terima Kasih}
Penulis menyampaikan terima kasih kepada Program Studi Sarjana Terapan Teknik Informatika, Universitas Logistik dan Bisnis Internasional (ULBI), serta seluruh pihak yang telah memfasilitasi pengujian infrastruktur jaringan pada penelitian Proyek 2 ini.

\bibliographystyle{IEEEtran}
\bibliography{references}

\end{document}
